
Compte tenu de la définition de @@M0@@, et de (6.6.4), on obtient :
@@M1@@
D'après (6.5.3), et le fait que @@M2@@, on obtient finalement :
@@M3@@
d'où le résultat.

\textbf{Proposition 6.10.} \uline{Soient} @@M4@@ \uline{un anneau binomial}, @@M5@@ \uline{une} @@M6@@-@@M7@@-\uline{algèbre augmentée}, @@M8@@ \uline{un élément de} @@M9@@ \uline{tel que} @@M10@@ \uline{pour} @@M11@@, \uline{et} @@M12@@. \uline{Alors} :
@@M13@@

On peut trouver des éléments @@M14@@, et @@M15@@ tels que :
@@M16@@
avec @@M17@@, et @@M18@@ pour tout @@M19@@ et tout @@M20@@. Soit @@M21@@ la @@M22@@-@@M23@@-algèbre engendrée par des générateurs @@M24@@, @@M25@@ soumis aux relations :
@@M26@@
On définit un homomorphisme d'augmentation @@M27@@ par @@M28@@. D'après (4.9.1), @@M29@@ est l'anneau des polynômes à coefficients dans @@M30@@ par rapport aux indéterminées @@M31@@, avec @@M32@@, et @@M33@@ avec @@M34@@, et la @@M35@@-filtration de @@M36@@ est identique à sa filtration par le poids, en attribuant le poids @@M37@@ à @@M38@@ et @@M39@@ (4.12 et 4.15).

On pose :
@@M40@@
